\section{A map of the reasons CBT fails}\label{sec:taxonomy}
``CBT does not work'' is not one claim. It bundles at least ten different ones, which differ in what they say fails (the average effect, the individual response, the mechanism, the evidence or the delivery), in how well they are supported, and in whether they are specific to CBT at all. Table~\ref{tab:reasons} lists them with the verdict this review reaches and the section that holds the evidence. We use four verdicts: \textbf{supported} (consistent evidence, quantified), \textbf{partly supported} (real but smaller or more conditional than the popular version), \textbf{contested} (credible evidence on both sides) and \textbf{open} (little or no direct evidence).

Two cautions apply throughout. First, many of the best-documented limits (dropout, relapse, comorbidity, harm under-reporting, therapist variability) are limits of psychological treatment in general, not of CBT in particular; we say so each time the source does. Second, a limit on certainty is not evidence of ineffectiveness: CBT keeps a moderate average effect under almost every sensitivity analysis we found \citep{plessen2023multiverse,carpenter2018cbt}. The map is therefore a list of \emph{where and why} CBT falls short of its reputation, not a verdict that it fails.

\begin{table}[ht]
\centering\small
\renewcommand{\arraystretch}{1.15}\begin{tabular}{@{}p{0.05\textwidth}>{\raggedright\arraybackslash}p{0.42\textwidth}>{\raggedright\arraybackslash}p{0.27\textwidth}>{\raggedright\arraybackslash}p{0.2\textwidth}@{}}
\toprule
 & Reason & Verdict & Where \\ \midrule
R1 & The cognitive skill demands executive function that anxiety (and ADHD) degrade & contested; direct moderator evidence open & Sections~\ref{sec:demands}--\ref{sec:adhd}, \ref{sec:moderators} \\
R2 & Cognitive restructuring is not the active ingredient & partly supported (underpowered trials) & Section~\ref{sec:ingredients} \\
R3 & Headline effects are inflated by waitlists, bias, low-quality trials and measurement choices & supported & Sections~\ref{sec:evidence}, \ref{sec:mechanism} \\
R4 & Many patients do not respond or drop out & supported (not CBT-specific) & Section~\ref{sec:nonresponse} \\
R5 & Gains erode after treatment ends & partly supported & Section~\ref{sec:nonresponse} \\
R6 & Harms and deterioration are under-recorded & supported for reporting; CBT lowers deterioration versus controls & Section~\ref{sec:nonresponse} \\
R7 & Patient factors (personality disorder, chronicity, adversity, medical comorbidity, neurodivergence) & partly supported; mostly not CBT-specific; autism evidence thin & Section~\ref{sec:patients} \\
R8 & Context: WEIRD trials, culture, waiting times, deprivation, access & supported & Section~\ref{sec:patients} \\
R9 & Delivery: therapist variability, drift, weak training, dose, format & partly supported; causal evidence missing & Section~\ref{sec:delivery} \\
R10 & The theory of change is under-tested; common factors may explain more & contested & Section~\ref{sec:mechanism} \\ \bottomrule
\end{tabular}
\vspace{4pt}\caption{Ten reasons given for why CBT fails, with the verdict of this review.}\label{tab:reasons}
\end{table}

The executive-function argument (R1) is the review's deepest case study, because it makes a testable, mechanism-level prediction and because it was the starting point of the project. It is one reason among ten, and Sections~\ref{sec:nonresponse}--\ref{sec:mechanism} treat the others at the same standard of evidence.
